\section{Model-based remedies and configuration: proofs}
\label{sec:annex-p16-model-based-remedy-and-configuration}

These arguments concern the source-typed, explicitly bounded constructions
in the reader, not unperformed correctness checks of all cited implementations.

\begin{annexitem}{annex:p16.rem-target-search}{thm:p16-rem-target-search}{All common targets, exact contexts and change-set minima}
\begin{proof}
The finite product has \(N\) assignments, including one empty assignment
when there are no reconfigurable components. Enumerate each at most once.
For common possibility, each target needs at most one consistency check
per hypothesis, giving \(kN\). For guarantee, test base consistency and
absence of a base model satisfying \(\neg G_i\), at most two per hypothesis.
Exactness makes these tests equivalent to nonempty base models all satisfying
the goal, yielding the \(2kN\) bound. Termination is conditional on the
total check premise, not inferred from finite component names alone.

A cached inconsistent partial state assignment is a conjunction contained
in every extending target. A witness of an acceptable target would satisfy
the matching background and that partial conjunction, contradicting its
inconsistency. In the guarantee case, the nonempty goal-entailing base has
a witness satisfying \(G_i\), so even a conflict found with goals added
is safe to use. A mere countermodel to entailment need not remain after
extra state constraints; it does not justify that partial-conflict pruning.
Context matching is essential because changed observations, goals or
structure can invalidate the cached inconsistency. Thus skipping removes
no acceptable common target; exact testing rejects every remaining
unacceptable one. Every acceptable target is consequently returned.

Checking the finite family for strictly smaller feasible change sets yields
exactly its inclusion-minimal members. Within the feasible subsets of a
cheapest target's change set, select a minimal feasible change set, which
exists by finiteness. Nondecreasing cost makes its cost no greater. This
uses the declared cost on the selected change-set representation; arbitrary
action-dependent downtime or irreversible risk is not proved monotone by it.
\end{proof}
\end{annexitem}

\begin{annexitem}{annex:p16.rem-plan-search}{thm:p16-rem-plan-search}{Every history prefix and the complete finite belief graph}
\begin{proof}
At length zero, \(B_0\) is exactly the admitted initial possibilities.
Inductively, a state in \(Post(B_j,a_j)\) has a predecessor with a
history witness for the prefix, and appending its transition supplies a
history of the longer prefix. Conversely every such history has a last
transition from a member of \(B_j\), so its endpoint belongs to the
same successor set. This proves both coverage directions, retaining the
model tags and nondeterministic alternatives.

Strong applicability gives at least one successor to every current
possibility; no initial world disappears merely because the action is
inapplicable. Inclusion of each full successor set in \(\mathcal S\)
therefore proves all-prefix safety. A nonempty final set contained in
\(\mathcal G\) makes every final history satisfy the terminal goal.
The initial safety test is necessary by the strong-plan definition.

There are at most \(2^n-1\) nonempty subsets of the finite tagged state
universe and at most \(m\) action trials from each, proving the graph
bounds and termination with visited-node reuse. Every strong plan follows
edges of this graph by its applicability/safety clauses. If a belief node
repeats along a plan, remove the cycle: the same remaining action sequence
has the same successor sets because these sets contain all relevant state
and history memory. Thus a successful graph path exists whenever a strong
plan exists. Conversely each stored success path is a strong plan by the
coverage/safety argument. Breadth-first expansion visits path lengths in
increasing order and finds a minimum-action-count success. No weighted
cost optimum follows from this unit-length ordering.

For the physical adequacy clause, represent the actual initial state in
\(B_0\). Inductively model applicability permits the real next action;
outcome coverage places its actual successor in the computed next set.
Decoded safety then holds after every prefix and decoded restoration at
the final state. Missing faults, wrong frames or an unrealizable action
break one of these premises, not the mathematical finite-graph proof.
\end{proof}
\end{annexitem}

\begin{annexitem}{annex:p16.rem-power-transfer}{thm:p16-rem-power-transfer}{Every first step and the four required actions}
\begin{proof}
Both endpoints have exactly one closed switch, so are safe. In the two-
switch case, the first action which changes a state either opens \(x\),
giving \((0,0)\) and losing critical supply, or closes \(y\), giving
\((1,1)\) and violating overload safety. No-op actions do not change
this obstruction. Thus no safe single-switch plan reaches the target.

With backup the displayed state sequence is
\[
 (1,0,0),(1,0,1),(0,0,1),(0,1,1),(0,1,0),
\]
in \((x,y,b)\) order. Every state has supply and at most one closed main route; every
backup guard holds by the declared healthy-available premise. It is a
safe target-reaching history. Any successful plan changes both \(x,y\)
at least once. Before the first main change, backup must be on, since
otherwise the same first-step obstruction holds. It must be off at the
end. At least two backup changes are therefore required in addition to
the two main changes, giving a four-action lower bound. Atomic simultaneous
main switching or a changed supply requirement is a different transition
model and is not excluded by this proof.
\end{proof}
\end{annexitem}

\begin{annexitem}{annex:p16.rem-policy-search}{thm:p16-rem-policy-search}{Ranks, every observation branch and the converse construction}
\begin{proof}
All nonempty safe beliefs are drawn from at most \(2^{|\mathcal X|}-1\)
subsets. Every nonfixed round adds at least one previously absent belief,
so only that many strict addition rounds can occur. At rank zero a belief
is safe and already terminal, and stopping is a depth-zero strong policy.
Inductively, a belief inserted at rank \(j+1\) has a strongly applicable
action with safe full successor set and every possible observation branch
of rank at most \(j\). Each successor belongs to exactly one observed
branch. The lower-rank policies thus cover all executions, maintain safety
and terminate within at most \(j+1\) actions.

Conversely a depth-zero strong policy must begin inside the goal and safety
sets, so lies in \(Win_0\). For a strong policy of depth at most \(j+1\),
its first action is strongly applicable and safe; every feasible observation
result has a continuation of depth at most \(j\). By induction every
branch belongs to \(Win_j\), so the starting belief belongs to
\(Win_{j+1}\). This proves both directions for each depth. Taking the
finite fixed point gives all and only beliefs with some finite-depth
strong policy. The argument imposes no positive probabilities or fairness;
it cannot discard a possible failing branch on those unprovided grounds.
\end{proof}
\end{annexitem}

\begin{annexitem}{annex:p16.rem-therapy}{thm:p16-rem-therapy}{All net-rate realizations and the immediate stock constraint}
\begin{proof}
Over the independent closed rate intervals, the largest value of \(p-k\)
is \(p_{max}-k_{min}\) and the smallest is \(p_{min}-k_{max}\),
with both corners admitted. All rates give a strictly negative derivative
exactly when the largest is negative; some do exactly when the smallest
is negative. This proves the two strict inequalities, including their
zero-boundary failures. The declared instantaneous persistence equation
gives \(q_1=q_0>H\), so immediate concentration below threshold is
inconsistent, irrespective of the new derivative. StopPump gives zero
supply, contradicting \(Q_1=Q_0>0\).

Even in the originally guaranteed-negative case, add an unmodeled positive
rate \(r>k_{min}-p_{max}\) at the maximal/minimal corner. The revised
derivative \(p_{max}-k_{min}+r\) is positive, proving the failure of
the old closure-dependent guarantee. A reservoir-located trigger is not
the tank-located instance generating its \(-k\) term; absent a declared
transport consequence the tank equation contains no such new influence.
Keeping the old \(q'_0>0\) as a time-one fact would also contradict
the intended negative \(q'_1\); the indexed before/after model avoids
that artificial conflict rather than retracting the actual old observation.
\end{proof}
\end{annexitem}

\begin{annexitem}{annex:p16.rem-configuration}{thm:p16-rem-configuration}{Every slot/count case and complete bounded construction}
\begin{proof}
With zero or one card, capacity is at most six, below ten. With two cards,
AA gives capacity eight and fails; BB gives power ten and fails; AB gives
capacity ten and power eight and succeeds. These are all two-card type
counts. The slots are identical and there is no order constraint, so AB
and BA are exactly the two ordered realizations. Removing either card
returns to a failing one-card case.

For the general bounded construction, take the finite product of all declared
choices, including empty-slot and connection/parameter choices. Each candidate
appears once. A total exact check terminates and accepts precisely the
candidates with a model of DD, that candidate and the goals. A sound
inconsistent partial design decision cannot be extended to a valid candidate,
since its model would also satisfy that inconsistent partial conjunction.
Therefore cached design conflicts safely prune, and every other candidate
is accepted or rejected exactly by its check. This proves finite termination
and both coverage directions for the bounded catalog, not for structures
whose permitted component count is unbounded or for absent functional laws.
\end{proof}
\end{annexitem}

\begin{annexitem}{annex:p16.rem-fmea}{thm:p16-rem-fmea}{Effect partition and all four redundant-route assignments}
\begin{proof}
A model of consistent \(\Gamma\) satisfies either E or its classical
negation. At least one of the two checks succeeds. If only the E check
succeeds, no model satisfies its negation and every model satisfies E;
if only the negated check succeeds, E is excluded; if both succeed,
there are witnesses for E and for its failure. These cases exhaust a
nonempty model class. Both failing instead indicates an inconsistent
fault/scenario model, not an effect guarantee.

In the route instance, the no-fault output is \((q_1,q_2,s)=(1,1,1)\).
Fault in route one gives \((0,1,1)\); fault in route two gives
\((1,0,1)\); faults in both give \((0,0,0)\). The deterministic
route laws, active-source scenario and OR law force exactly these tuples,
so each is a nonempty unique observable realization. E, which is \(s=0\),
is excluded in the first three and guaranteed in the fourth. Projection
to \(s\) gives the same singleton \(\{1\}\) for normal and each
single-fault hypothesis, so that observation cannot refute either single
in favor of normal. Projection to the two route outputs yields four
different tuples and distinguishes all assignments. Singles-only
enumeration therefore omits the exhibited double-fault loss of service.
\end{proof}
\end{annexitem}

\begin{annexitem}{annex:p16.rem-debugging}{thm:p16-rem-debugging}{The missing increment and the finite-test countermodel}
\begin{proof}
Combining \(y=x\) with \(y=x+1\) would imply \(0=1\), impossible
over the stated integers for any shared input. In the catalog, the
input-zero output is \(\delta\); expecting one excludes zero and
accepts one. For \(\delta=1\), the program relation is exactly
\(y=x+1\) for every integer, proving the full claimed specification
at that catalog's scope. The overfitting program has output one at zero,
so satisfies the same finite expectation, but at input one returns one
instead of two. This is a complete countermodel to identifying a passing
test with general correctness. The old program's observed zero output
is a fact about its old version and cannot be imposed as the new version's
output without changing the repair task back into an artificial conflict.
\end{proof}
\end{annexitem}
